\section{Feature List：结构化的任务追踪}

Feature list 是整个方案中最关键的组件之一，它直接解决了 Agent 过早宣布完工和试图一次性完成所有任务的问题。

\subsection{设计理念}

Initializer Agent 会基于用户的初始 prompt，生成一份全面的功能需求清单。在"克隆 claude.ai"的示例中，这份清单包含了超过 200 个功能点，例如"用户可以打开新对话，输入问题，按回车，看到 AI 回复"。

每个功能点最初都被标记为 \texttt{"passes": false}，这样后续的 Coding Agent 就能清楚地看到完整功能的全貌，而不会因为部分功能已经实现就误以为项目完成。

\subsection{JSON 格式的选择}

\begin{lstlisting}[language={}, caption={Feature list 中单个功能的 JSON 结构}]
{
  "category": "functional",
  "description": "New chat button creates a fresh conversation",
  "steps": [
    "Navigate to main interface",
    "Click the 'New Chat' button",
    "Verify a new conversation is created",
    "Check that chat area shows welcome state",
    "Verify conversation appears in sidebar"
  ],
  "passes": false
}
\end{lstlisting}

\begin{importantbox}{为什么选择 JSON 而非 Markdown？}
Anthropic 团队经过实验发现，与 Markdown 文件相比，模型不太可能不恰当地修改或覆盖 JSON 文件的内容。JSON 的结构化特性使得 Agent 更倾向于只修改 \texttt{passes} 字段的值，而不是重写整个描述。这是一个值得借鉴的 prompt engineering 技巧：\textbf{数据格式的选择会影响模型的行为倾向}。
\end{importantbox}

\subsection{严格的修改规则}

Coding Agent 被明确指示只能修改 \texttt{passes} 字段的值——从 \texttt{false} 改为 \texttt{true}。Anthropic 团队在 prompt 中使用了强措辞：

\begin{quote}
\textit{"It is unacceptable to remove or edit tests because this could lead to missing or buggy functionality."}
\end{quote}

这种措辞上的强调对于控制 Agent 行为非常有效。它明确告诉 Agent，删除或修改功能描述是不可接受的，从根本上防止了 Agent 通过降低标准来"完成"任务的取巧行为。

\begin{warningbox}{Agent 的"作弊"倾向}
在没有明确约束的情况下，Agent 可能会通过删除或降低测试标准来更快地"完成"任务。这类似于人类中的 Goodhart 定律——当指标本身成为目标时，它就不再是一个好的指标。在 Agent 系统设计中，必须通过严格的规则和不可变的评判标准来防止这种行为。
\end{warningbox}

\subsection{本章小结}

Feature list 通过结构化的 JSON 格式将项目需求具体化为可追踪的功能点，每个功能都有清晰的验证步骤。JSON 格式的选择、初始状态全部标记为失败、以及严格禁止修改功能描述的规则，共同构成了防止 Agent 过早宣布完工的有效机制。


